\section{Lista de riesgos y biblioteca de antipatrones: evitar convertir un Agente de propósito general en un «sistema de demostración»}

\subsection{Antipatrones a nivel de producto}

La entrevista advierte repetidamente sobre un riesgo: muchos equipos confunden «poder hacer una demostración» con «poder entregar». Lo primero busca un efecto impactante a corto plazo; lo segundo busca estabilidad sostenida a largo plazo. Los antipatrones más comunes a nivel de producto incluyen: optimizar solo la primera respuesta, no gestionar el estado de las tareas largas, no exponer las causas de las fallas y no proporcionar resultados intermedios que se puedan revisar. A corto plazo estas prácticas reducen la dificultad del desarrollo, pero a largo plazo hacen que la confianza del usuario disminuya de manera continua.

\begin{warningbox}{4 grandes antipatrones a nivel de producto}
\begin{enumerate}
    \item Demo-first: toda la optimización gira en torno a escenarios de exhibición, no a la distribución real de las tareas
    \item Happy-path-only: solo se cubre el camino sin contratiempos, sin manejo de excepciones
    \item Opaque-output: la salida parece completa, pero carece de evidencia rastreable
    \item Metric-theater: las métricas se ven bien, pero la retención y la recompra de los usuarios bajan en lugar de subir
\end{enumerate}
\end{warningbox}

\subsection{Antipatrones a nivel técnico}

Los antipatrones a nivel técnico provienen principalmente de una «complejidad fuera de control»: cada vez se agregan más herramientas y la política se vuelve cada vez más difícil de aprender; la capa de orquestación se vuelve cada vez más gruesa y la localización de fallas cada vez más lenta; cuando la capacidad del modelo se actualiza, nadie gestiona la deuda de la capa del sistema. En la entrevista, acciones aparentemente conservadoras como «retirar herramientas», «hacer repeticiones» y «enfatizar la observabilidad» son en realidad una gestión activa de la complejidad.

\begin{knowledgebox}{Principios para gobernar la complejidad a nivel técnico}
\begin{itemize}
    \item El camino predeterminado debe ser corto: primero hay que asegurar que un camino sea estable, y después expandir las ramificaciones
    \item Las fallas deben ser localizables: cada falla debe poder ubicarse en un paso concreto y en un límite de responsabilidad
    \item Los cambios deben poder revertirse: cualquier modificación de la política debe poder deshacerse con rapidez
\end{itemize}
\end{knowledgebox}

\subsection{Antipatrones a nivel organizacional}

Los riesgos a nivel organizacional suelen manifestarse como «divergencia de objetivos + desequilibrio de ritmo». Cuando un equipo persigue demasiadas direcciones a la vez, las prioridades existen solo de nombre y en la práctica desaparecen; cuando falta un mecanismo de revisión posterior, las fallas reaparecen de distintas formas. El marco GPA y la «lista de cosas que no se hacen» que aparecen en la entrevista resuelven, en esencia, el problema del aumento de la entropía a nivel organizacional.

\begin{importantbox}{Acciones para prevenir la deriva a nivel organizacional}
\begin{itemize}
    \item Cada semana se conservan solo unos pocos objetivos centrales; el resto de los asuntos se pone explícitamente en espera
    \item La revisión posterior debe producir «elementos que se detienen la próxima semana», no solo «elementos planeados para la próxima semana»
    \item Los puestos clave comparten el mismo panel de indicadores, para evitar óptimos locales
\end{itemize}
\end{importantbox}

\subsection{Matriz de prioridad de riesgos}

\begin{center}
\begin{tabular}{p{3.2cm}p{2.1cm}p{2.1cm}p{5.4cm}}
\toprule
Elemento de riesgo & Probabilidad de ocurrencia & Nivel de impacto & Acción prioritaria de gobierno \\
\midrule
Falla silenciosa en tareas largas & Alta & Alta & Agregar monitoreo de latidos, alertas por tiempo de espera y capturas automáticas \\
Invocación incorrecta de herramientas & Alta & Media & Establecer una lista blanca de herramientas, un camino predeterminado y auditoría de invocaciones \\
Salida no revisable & Media & Alta & Exigir que se adjunten resultados intermedios y referencias de evidencia \\
La organización combate en varios frentes & Media & Alta & Ejecutar la lista de cosas que no se hacen, reducir el número de líneas principales \\
Divergencia entre métricas y experiencia & Media & Media & Introducir revisión humana y un panel de calidad subjetiva \\
Regresiones de versión invisibles & Media & Alta & Establecer un conjunto de tareas de regresión y una puerta de calidad para el lanzamiento \\
\bottomrule
\end{tabular}
\end{center}

\subsection{Resumen de la sección}

Hacer explícita la lista de riesgos es el parteaguas que lleva a un equipo de Agentes de propósito general de «saber construir funciones» a «saber construir sistemas». El verdadero valor que aporta la entrevista es permitir que el equipo cuente con la capacidad de gobernar la complejidad antes de expandirse.

